\infographicchapter{Corneal Tomography}{corneal_tomography.png}
\label{ch:infographic-corneal-tomography}

\section{Overview}

Corneal tomography is a quick, completely non-touch 3D eye scan that maps both the front and back surfaces of your cornea (the clear, dome-shaped window at the front of your eye, which provides roughly 70\% of the eye's total focusing power) along with its full thickness from edge to edge. Modern scanners work in one of two ways: a rotating Scheimpflug camera tilts its lens to keep the entire curved front of the eye in crisp focus as a gentle slit of light sweeps across it, or an anterior segment optical coherence tomography (AS-OCT) scanner uses harmless near-infrared light echoes to measure corneal layers with microscopic accuracy.

Unlike simpler scans that look only at the front curve of the eye, corneal tomography shows bumps or dips on both the front and back surfaces, maps where the cornea is thickest and thinnest, and calculates its total focusing power. Its most important jobs are catching corneal thinning diseases early---especially keratoconus, where subtle bulging on the back surface and localized thinning often show up before the front surface changes---and screening patients before laser vision surgery (such as LASIK, PRK, or SMILE) to make sure surgery is safe. Doctors also rely on it to watch whether keratoconus is worsening (to time corneal cross-linking treatment), check corneal transplants, and fit custom contact lenses.

The biggest benefits are that it is fast, never touches your eye, uses no X-ray radiation, creates a complete 3D map, and includes built-in screening scores that catch early warning signs. Its main limits are that blinking, dry eyes, eyelid shadows, or recent contact lens wear can skew the maps, and thick corneal scars can scatter light (though infrared AS-OCT sees through haze better than visible-light cameras). Because it uses only safe visible or near-infrared light, corneal tomography carries virtually no procedural risk.

\section{Underlying Physical Principles}
\begin{itemize}
    \item \textbf{Tilted-Lens Photography (Scheimpflug Systems):} Using the Scheimpflug geometric principle (named after Theodor Scheimpflug), a rotating camera tilts its lens and sensor planes so they intersect the plane of focus along a shared line rather than sitting parallel. This tilted alignment keeps an entire angled cross-section of the eye---from the outer cornea back to the natural lens---sharply in focus all at once as a thin slit of blue light sweeps across it.
    \item \textbf{Light-Echo Measurement (AS-OCT Systems):} Anterior segment OCT scanners shine harmless near-infrared light (840--1,310 nm) into the eye and measure tiny timing differences in the light bouncing back from the front and back corneal surfaces, achieving 5--15 micrometer accuracy even through mild cloudiness.
    \item \textbf{Accounting for Light Bending:} As light crosses the curved front of the cornea, it bends (refracts). Built-in optical formulas correct for that bending so the computer pinpoints the exact true shape of the back surface.
    \item \textbf{3D Sweep Around the Clock:} Over 2--5 seconds, the scanner captures 25--100 cross-sectional slices arranged like spokes of a wheel around the center of the cornea, then compares both surfaces to an ideal smooth reference dome to highlight subtle high or low spots.
    \item \textbf{Thickness and Total Focusing Power:} Subtracting the back surface from the front surface at every point creates a full-cornea thickness map (pachymetry), while combining both curves calculates the cornea's true total focusing power.
\end{itemize}

\section{Imaging Mechanism}
\begin{enumerate}
    \item \textbf{Easy Non-Touch Alignment:} You rest your chin and forehead on the supports while an automatic eye-tracking camera lines up gently with the center of your eye as you look at a small light.
    \item \textbf{Quick Rotational Sweep:} A rotating camera sweeps a thin blue slit of light across your cornea over 2--5 seconds (or an AS-OCT scanner sweeps an invisible near-infrared light beam across the front of your eye).
    \item \textbf{Recording Cross-Sections:} High-speed digital sensors capture sharp slice images at every angle while built-in software checks automatically for blinks or eye movement.
    \item \textbf{Building the 3D Cornea:} The computer outlines the front and back corneal edges, adjusts for light bending, aligns all slices into a single 3D model, and calculates thickness and surface elevation maps.
    \item \textbf{Color-Coded Map Display:} The screen shows easy-to-read color maps of front curvature, front elevation, back elevation, and thickness side by side, along with automated keratoconus screening scores (such as the Belin--Ambrosio display).
\end{enumerate}

\section{Key Technical Parameters}
\noindent
\begin{tabularx}{\textwidth}{@{}l X@{}}
\toprule
\textbf{Parameter} & \textbf{Typical Value / Description} \\
\midrule
Physical Principle & Rotating tilted-lens slit photography (Scheimpflug) or near-infrared light-echo imaging (AS-OCT) \\
Energy / Signal Source & Safe blue LED slit light (Scheimpflug) or low-power near-infrared light (AS-OCT) \\
Typical Frequency / Energy & 475 nm visible blue light (Scheimpflug) or 840--1,310 nm near-infrared light (AS-OCT) \\
Spatial Resolution & About 5--15 micrometers depth accuracy (AS-OCT) to 20--50 micrometers elevation accuracy (Scheimpflug) \\
Typical Scan Duration & 2--5 seconds per eye (5--10 minutes total for both eyes) \\
Contrast Agent & None (completely non-touch light scan) \\
Ionizing Radiation & None \\
\bottomrule
\end{tabularx}

\section{Common Symptoms That Prompt This Scan}
\begin{itemize}
    \item \textbf{Progressive Astigmatism, Glare, or Double-Edged Vision:} Ghosting around lights, nighttime halos, or rapidly shifting prescriptions needing a full front-and-back 3D corneal map to catch early keratoconus.
    \item \textbf{Pre-Surgery Safety Screening (No Symptoms):} Essential before LASIK, SMILE, PRK, or premium cataract surgery to confirm the cornea is thick and symmetric enough for safe treatment.
    \item \textbf{Monitoring Known Corneal Thinning:} Tracking whether corneal bulging is stable or progressing so sight-saving cross-linking treatment can be timed properly.
    \item \textbf{Cloudy or Swollen Cornea:} Morning blurry vision or haze when doctors need precise full-thickness corneal maps in Fuchs dystrophy or after corneal transplant.
    \item \textbf{Family History, Teens, or Chronic Eye Rubbing (Often No Symptoms):} Screening children, teenagers, and young adults who have a close relative with keratoconus (about 7.7\% of first-degree relatives are affected, compared with roughly 54.5 per 100,000 people in the general US population) or who have eye allergies and frequent eye rubbing, which can trigger or speed up corneal thinning.
\end{itemize}

\section{Clinical Usages}

\subsection{Primary Usages}
\begin{itemize}
    \item \textbf{Catching Keratoconus Early:} Spots early, hidden (subclinical) and established keratoconus, pellucid marginal degeneration, and post-surgery corneal bulging by detecting subtle back-surface protrusion and localized thinning---which is especially vital in children and teenagers, where thinning progresses faster and timely cross-linking can help avoid a corneal transplant.
    \item \textbf{Safety Screening Before Laser Eye Surgery:} Mandatory 3D check before LASIK, PRK, or SMILE to rule out hidden corneal weakness and make sure enough healthy corneal thickness will remain after surgery.
    \item \textbf{Tracking Disease Progression:} Compares scans over time to see if the cornea is getting steeper or thinner, helping doctors decide the best moment to perform corneal cross-linking (a procedure that strengthens the cornea).
    \item \textbf{Planning Cataract and Specialty Lens Surgery:} Measures true total corneal focusing power---including astigmatism on the back surface---to select the most accurate toric or multifocal lens implants during cataract surgery.
    \item \textbf{Full-Cornea Thickness Mapping:} Maps thickness across the entire cornea to evaluate corneal swelling (such as in Fuchs dystrophy) and interpret eye-pressure readings accurately in glaucoma.
\end{itemize}

\subsection{Secondary Usages}
\begin{itemize}
    \item \textbf{Checking the Front Chamber and Drainage Angle:} Measures the depth and drainage angle of the front of the eye to screen for narrow-angle glaucoma and size implantable contact lenses.
    \item \textbf{Corneal Transplant Planning and Care:} Helps surgeons match donor tissue thickness and depth for partial-thickness transplants and monitor healing afterward.
    \item \textbf{Planning Corneal Ring Implants:} Measures exact regional thickness to guide safe placement of tiny supportive ring segments in keratoconus.
    \item \textbf{Custom Contact Lens Fitting:} Supplies 3D height and shape data to design comfortable scleral, hybrid, and rigid contact lenses for irregular or transplanted corneas.
\end{itemize}

\section{Advantages}
\begin{itemize}
    \item \textbf{Complete 3D Picture:} Maps the front surface, back surface, and full thickness of the cornea all at once, rather than just the front curve.
    \item \textbf{Spots Hidden Weakness Early:} Catches subtle back-surface bulging and localized thinning before the front of the eye changes, making refractive surgery screening much safer.
    \item \textbf{Non-Touch, Fast, and Completely Safe:} Uses harmless visible or near-infrared light with zero eye contact and zero X-rays, so it can be repeated comfortably as often as needed.
    \item \textbf{Smart Built-In Screening Scores:} Validated screening displays (like Belin--Ambrosio indices) help doctors spot subtle keratoconus and track changes consistently over time.
    \item \textbf{Sharper Cataract Lens Calculations:} Accounts for back-surface astigmatism, improving vision outcomes with specialty cataract lens implants---especially in eyes that had past laser surgery.
\end{itemize}

\section{Disadvantages}
\begin{itemize}
    \item \textbf{Sensitive to Dry Eyes and Contact Lenses:} An uneven tear film, severe dry eye, mucus, or recent contact lens wear can distort both front and back measurements.
    \item \textbf{Cloudy Scars Scatter Light:} Dense corneal scars, haze, or swelling scatter visible light and make back-surface tracking harder, especially on Scheimpflug cameras compared with infrared AS-OCT.
    \item \textbf{Blinking and Eyelid Shadows:} Looking away, blinking, tight eyelids, or long lashes can block the edges of the scan and create false high spots.
    \item \textbf{Machines Are Not Directly Interchangeable:} Numbers and screening scores can differ slightly between brands or between Scheimpflug and AS-OCT machines, so repeat checkups should be done on the same type of scanner.
    \item \textbf{Higher Cost:} Tomographers cost more than simpler front-surface topographers and require careful quality checks to avoid false alarms.
\end{itemize}

\section{Contraindications}

\subsection{Absolute}
\begin{itemize}
    \item \textbf{None:} Because corneal tomography never touches your eye, uses low-energy harmless light, and needs no radiation or dye, it has no absolute medical contraindications.
\end{itemize}

\subsection{Relative}
\begin{itemize}
    \item \textbf{Contagious Eye Infection:} Non-urgent scans should wait if you have a contagious eye infection (like viral pink eye) so germs are not spread through the chinrest or equipment.
    \item \textbf{Recent Contact Lens Wear:} Scanning for laser eye surgery without a proper break from contact lenses (usually 1--2 weeks for soft lenses, 2--4 weeks for rigid lenses) is strongly discouraged because lenses temporarily mold the cornea and can hide or imitate disease.
    \item \textbf{Severe Light Sensitivity or Trouble Sitting Still:} Extreme light sensitivity, rapid involuntary eye shaking, or inability to sit at a chinrest can prevent a clear 3D scan.
\end{itemize}

\section{Patient Safety Guidelines}
\begin{itemize}
    \item \textbf{Safe Light Levels:} Scheimpflug and AS-OCT scanners stay well within FDA and ANSI eye light-safety limits with zero X-rays, though technicians keep repeat flashes minimal for very light-sensitive or freshly operated eyes.
    \item \textbf{Cleanliness Between Patients:} Following AAO hygiene standards, staff wipe down the chinrest, forehead bar, and handles with approved disinfectant wipes between every patient and reschedule non-urgent scans during contagious pink eye.
    \item \textbf{Contact Lens Break for Surgical Safety:} Enforcing a full lens-free break---typically 1--2 weeks for soft lenses and 2--4 weeks for rigid lenses---prevents lens molding from masking hidden keratoconus before laser eye surgery.
    \item \textbf{Gentle, Non-Touch Technique:} Technicians keep the camera safely back from the eye, avoid pressing on the eyeball if lifting a droopy lid, and treat severe dry eye before taking critical surgical measurements.
    \item \textbf{Quality Checks and Same-Machine Follow-Up:} Clinicians verify built-in scan quality flags before making decisions and use the same scanner brand for follow-up visits so real changes are not confused with machine differences.
\end{itemize}

\section{Prerequisite Tests}
\begin{itemize}
    \item \textbf{No Blood Tests Needed:} Because this is a non-touch light scan with no radiation or dye, no blood work, kidney checks, or pregnancy tests are needed.
    \item \textbf{Standard Eye Exam First:} Per AAO refractive surgery guidelines, your doctor first checks your vision, glasses prescription, and front-of-eye health under a slit-lamp microscope to assess corneal clarity and tear quality.
    \item \textbf{Confirming Your Contact Lens Break:} Before laser surgery screening or keratoconus testing, staff confirm you have kept contact lenses out long enough (1--2 weeks for soft, 2--4 weeks or longer for rigid) for your cornea to relax back to its natural shape.
    \item \textbf{Checking for Dry Eye:} Evaluating and calming significant dry eye or eyelid irritation beforehand helps prevent false bumps or thin spots on the map.
\end{itemize}

\section{Preparing for the Scan}
\begin{itemize}
    \item \textbf{Take a Break From Contact Lenses:} Follow your clinic's lens-free window closely---typically 1--2 weeks for soft lenses and 2--4 weeks (or until repeat maps stop changing) for rigid or scleral lenses---so temporary lens molding does not skew your results.
    \item \textbf{No Fasting or Medicine Changes:} Eat, drink, and take your regular medications as usual.
    \item \textbf{Skip Eye Makeup and Thick Ointments:} Avoid mascara, eyeliner, or heavy eye creams that can flake into your tears. You may continue prescribed preservative-free artificial tears, but avoid thick eye gels or ointments right before the scan.
    \item \textbf{Easy Positioning:} You will rest your chin and forehead steadily on the clean supports, blink naturally right before the scan to smooth your tear film, then open wide and look steadily at the target light for a few seconds.
    \item \textbf{No Driver Needed:} Standalone corneal tomography uses no dilating or numbing drops, so you can drive home right afterward (just keep contact lenses out if your clinic asked you to wait for surgical clearance or a repeat check).
\end{itemize}

\section{During Scan}
\begin{itemize}
    \item \textbf{Chinrest and Target Light:} You sit comfortably with your chin on the rest and forehead lightly against the bar, opening both eyes wide while looking steadily at the glowing target inside the machine.
    \item \textbf{Easy Breathing and Timed Blinks:} Breathe normally. The technician will ask you to blink once or twice to refresh your tears, then keep your eyes wide open without blinking or squinting for just a few seconds.
    \item \textbf{What You Will See and Hear:} Nothing touches your eye and the test is completely painless. You will see a bright rotating blue or white bar of light (on a Scheimpflug scanner) or a gentle red or invisible fixation light (on an AS-OCT scanner) along with a soft whirring sound.
    \item \textbf{Staying Comfortable:} Tell the technician if the bright light makes your eyes water or if the table height strains your neck, so they can adjust your chair and give you a quick breather between eyes.
    \item \textbf{How Long It Takes:} Scanning each eye takes just 2--5 seconds, and the whole visit for both eyes usually wraps up in 5--10 minutes.
\end{itemize}

\section{Normal Appearance}
\begin{itemize}
    \item \textbf{Smooth Front and Back Elevation:} Both surfaces stay smooth and symmetrical near the ideal reference curve (usually within about +12 micrometers on the front and +15--18 micrometers on the back on standard Scheimpflug maps), with no isolated bulging ``islands.''
    \item \textbf{Even Thickness Rings:} Thickness forms neat, centered rings---typically 520--560 micrometers near the center---with the thinnest spot close to the middle and smooth, gradual thickening out toward the edges.
    \item \textbf{Regular Front Curvature:} Shows an even round, oval, or symmetrical bowtie shape (for normal astigmatism) with balanced top-and-bottom focusing power (roughly 41--46 diopters).
    \item \textbf{Green-Light Screening Scores:} Composite ectasia scores (such as a Belin--Ambrosio final ``D'' score below 1.6) fall cleanly in the normal range with no yellow or red warning flags.
\end{itemize}

\section{Observing for Abnormal Patterns}
\begin{itemize}
    \item \textbf{Back-Surface Bulge (Early Keratoconus):} A small, off-center ``island'' bulging outward on the back-surface map (usually downward and slightly toward the ear) is one of the earliest clues to hidden or developing keratoconus, often appearing before the front curve changes.
    \item \textbf{Off-Center Thinning:} The thinnest spot shifting downward and outward, a thinner-than-normal cornea, a rapid jump in thickness from center to edge, and elevated Belin--Ambrosio ``D'' scores strongly point to corneal thinning (ectasia).
    \item \textbf{Pellucid Marginal Degeneration:} A narrow curved band of thinning near the very bottom edge of the cornea (1--2 mm above the lower rim) paired with a classic ``crab-claw'' curvature shape sets this condition apart from central keratoconus.
    \item \textbf{Early Corneal Swelling (Fuchs Dystrophy):} Uneven thickness rings, the thinnest spot shifting out of its usual position, and gentle backward bulging of the inner surface can reveal early corneal swelling before visible blisters appear.
    \item \textbf{Harmless Look-Alikes:} Recent contact lens wear can temporarily press and mold the cornea into an uneven shape that looks just like keratoconus. Dry eye, eyelid shadows, and dense scars can also trick the camera into showing false thin spots or back-surface bumps.
\end{itemize}

\section{Dealing with Uncertainties}
\begin{itemize}
    \item \textbf{Dry Eye, Shadows, and Scars:} Broken tear film, mucus, droopy lids, long lashes, off-center gazing, or corneal scars can scatter light and create false elevation bumps or falsely thin readings.
    \item \textbf{Contact Lens Molding vs.\ Early Disease:} Contact lens warpage can closely imitate keratoconus, while very early keratoconus can sit right on the borderline of normal on a single chart.
    \item \textbf{Different Brands Give Slightly Different Numbers:} Comparing a scan from one brand to a different machine can make a stable eye look like it changed when only the hardware differed.
    \item \textbf{Double-Checking Quality:} Doctors resolve borderline scans by checking quality indicators, calming dry eye, extending the contact lens break (2--4 weeks or longer), and repeating scans on the exact same machine until results match consistently.
    \item \textbf{Combining Tests and Watching Safely:} For borderline laser-surgery candidates or uncertain maps, doctors add front-surface topography, AS-OCT surface-cell (epithelial) thickness maps, and corneal stiffness tests. If doubt remains, laser surgery is safely postponed and the eye is rechecked every 3--6 months.
\end{itemize}

\section{Augmenting Methods}
\begin{itemize}
    \item \textbf{AS-OCT Surface-Layer (Epithelial) Mapping:} Measures the cornea's ultra-thin outer skin layer, which naturally thins over an early hidden cone (forming a ``doughnut'' ring) to unmask early keratoconus even sooner.
    \item \textbf{Placido Disc Corneal Topography:} Projects concentric light rings to give extra-sensitive front-surface curvature and tear-film smoothness checks that cross-verify tomographic maps.
    \item \textbf{Corneal Stiffness (Biomechanical) Testing:} Uses a gentle puff of air to measure how flexible or stiff the cornea is, adding strength information to the 3D shape map when screening for laser surgery.
    \item \textbf{Specular Microscopy:} Counts and inspects the cornea's innermost pump cells alongside thickness maps when checking Fuchs dystrophy or planning cataract surgery.
    \item \textbf{Smart Screening Software:} Multivariable and AI models combine front shape, back shape, thickness, and stiffness to spot early keratoconus with higher accuracy.
\end{itemize}

\section{Comparison with Alternative Modalities}
\begin{itemize}
    \item \textbf{Versus Placido Disc Corneal Topography:} Tomography maps the front surface, back surface, and full thickness---making it essential for laser-surgery screening and early keratoconus detection. Placido topography maps only the front surface, though it is very sensitive to subtle surface and tear-film ripples.
    \item \textbf{Versus Ultrasound Pachymetry:} Ultrasound thickness testing requires numbing drops, touches the eye, and measures only one spot at a time, whereas tomography never touches the eye and maps thickness across the entire cornea to find the true thinnest point.
    \item \textbf{Versus Standalone AS-OCT:} Scheimpflug tomography offers long-established keratoconus screening scores, while infrared AS-OCT sees better through cloudy scars and adds microscopic outer-layer thickness mapping.
\end{itemize}

\section{Follow-up Steps}
\begin{itemize}
    \item \textbf{Right After the Scan:} Because nothing touched your eye and no drops were used, you can return to normal activities immediately (keeping contact lenses out only if your clinic asked you to stay lens-free for upcoming surgery or repeat scans).
    \item \textbf{Laser Surgery and Keratoconus Decisions:} Your doctor reviews the 3D maps alongside your prescription, eye exam, and risk factors before clearing you for LASIK, PRK, or SMILE. Borderline or abnormal maps mean laser surgery is postponed and a cornea specialist evaluates you.
    \item \textbf{Companion Tests When Needed:} Doctors may add Placido topography, AS-OCT epithelial mapping, or corneal stiffness tests for borderline thinning, or specular microscopy if Fuchs dystrophy is suspected.
    \item \textbf{Scheduled Checkups, Eye-Rubbing Control, and Family Screening:} Younger patients or those with active keratoconus usually repeat tomography on the same machine every 3--6 months (and stable adults every 6--12 months) so corneal cross-linking can be done promptly if any progression appears. Your doctor may also recommend treating eye allergies to stop chronic eye rubbing and checking close family members for hidden corneal thinning.
\end{itemize}

\section{Key References}
\begin{itemize}
    \item American Academy of Ophthalmology (AAO). \textit{Preferred Practice Pattern: Refractive Errors \& Refractive Surgery} and \textit{Preferred Practice Pattern: Corneal Ectasia}.
    \item Gomes JAP, Tan D, Rapuano CJ, et al. \textit{Global Consensus on Keratoconus and Ectatic Diseases}. Cornea.
    \item Belin MW, Ambrosio R Jr. \textit{Scheimpflug Imaging for Keratoconus and Ectatic Disease}. Indian J Ophthalmol; and StatPearls Publishing: \textit{Keratoconus}.
    \item American National Standards Institute (ANSI). \textit{ANSI Z80.23: Ophthalmic Instruments---Corneal Topography Systems}.
    \item Magklaras E, Karamitsou K, Diakonis VF, Mprotsis T, Tsaousis KT. \textit{Early Detection of Keratoconus: Diagnostic Advances and Their Impact on Visual Outcomes: A Systematic Review}. Medicina (Kaunas) / National Library of Medicine (NIH PMC). Available at: \url{https://pmc.ncbi.nlm.nih.gov/articles/PMC12843096/}.
\end{itemize}

\section{Further Reading}
\begin{itemize}
  \item Wikipedia contributors. \textit{Corneal topography}. Wikipedia, The Free Encyclopedia. Available at: \url{https://en.wikipedia.org/wiki/Corneal_topography}.
  \item Wikipedia contributors. \textit{Scheimpflug principle}. Wikipedia, The Free Encyclopedia. Available at: \url{https://en.wikipedia.org/wiki/Scheimpflug_principle}.
\end{itemize}
